\documentclass[letterpaper,11pt]{article}
\usepackage[T1]{fontenc}
\usepackage[utf8]{inputenc}
\usepackage{lmodern}
\usepackage[french]{babel}
\usepackage[margin=1in,headheight=28pt]{geometry}
\renewcommand{\familydefault}{\sfdefault}
\usepackage{amsmath,amssymb,amsthm,bm,enumitem}
\usepackage[dvipsnames]{xcolor}
\usepackage[small,compact]{titlesec}
\usepackage{fancyhdr,fancyvrb}
\usepackage[colorlinks=true,allcolors=blue]{hyperref}
\pagestyle{fancy}
\fancyhead[L]{STT-2100 \\ Régression}
\fancyhead[R]{A.-A. Vallée \\ Université Laval}
\fancyfoot[C]{\thepage}
\setlength{\parindent}{0pt}
\setlength{\parskip}{5pt}
\fvset{fontsize=\small,frame=single,framesep=2mm}
\newcommand{\esp}{\mathbb{E}}
\newcommand{\var}{\operatorname{Var}}
\newcommand{\bo}[1]{\bm{#1}}
\newcommand{\N}{\mathcal{N}}
% Définir \avecsolutions avant l'input pour une version corrigée.
\newcommand{\prof}[1]{\ifdefined\avecsolutions\par{\color{blue}\textbf{Solution :} #1}\par\fi}
\begin{document}
\begin{center}{\Large Diagnostiquer la multicolinéarité avec R}\end{center}
Cette fiche accompagne les notes du chapitre 9. Les noms des objets peuvent être adaptés à vos données. Les exemples utilisent des données incluses dans R ou des données simulées; aucun fichier externe n’est requis.

\par\medskip\noindent\begin{minipage}{\linewidth}
\subsection*{Dépendance exacte et rang}
Une corrélation par paire ne suffit pas à dépister les dépendances entre plusieurs colonnes. alias() et le rang QR révèlent une dépendance exacte.
\begin{Verbatim}
donnees <- mtcars
donnees$double_wt <- 2 * donnees$wt
mod_exact <- lm(mpg ~ wt + double_wt + hp, data = donnees)
coef(mod_exact)
alias(mod_exact)
X <- model.matrix(mod_exact)
c(rang = qr(X)$rank, colonnes = ncol(X))
\end{Verbatim}

\end{minipage}\par
\par\medskip\noindent\begin{minipage}{\linewidth}
\subsection*{VIF par régressions auxiliaires}
Cette formule concerne des prédicteurs numériques à une colonne, dans un modèle avec constante, sans interaction et de plein rang. Pour un facteur à plusieurs colonnes, un diagnostic par terme tel que le GVIF exige une autre formule.
\begin{Verbatim}
variables <- c("wt", "hp", "disp")
cor(donnees[variables])
vif <- sapply(variables, function(v) {
  f <- reformulate(setdiff(variables, v), response = v)
  auxiliaire <- lm(f, data = donnees)
  1 / (1 - summary(auxiliaire)$r.squared)
})
vif
\end{Verbatim}

\end{minipage}\par
\par\medskip\noindent\begin{minipage}{\linewidth}
\subsection*{Indices de conditionnement}
On normalise ici toutes les colonnes, constante comprise, à une norme euclidienne de 1. Les indices sont les racines des rapports de valeurs propres. Le maximum est le conditionnement spectral de cette matrice de modèle normalisée.
\begin{Verbatim}
mod <- lm(mpg ~ wt + hp + disp, data = donnees)
X <- model.matrix(mod)
Xs <- sweep(X, 2, sqrt(colSums(X^2)), "/")
decomp <- eigen(crossprod(Xs), symmetric = TRUE)
if (min(decomp$values) <= 0) stop("Matrice singuliere")
indices <- sqrt(max(decomp$values) / decomp$values)
contributions <- sweep(decomp$vectors^2, 2,
                       decomp$values, "/")
proportions <- contributions / rowSums(contributions)
rownames(proportions) <- colnames(X)
indices
proportions
\end{Verbatim}

\end{minipage}\par
\par\medskip\noindent\begin{minipage}{\linewidth}
\subsection*{Une pénalisation ridge illustrative}
Les prédicteurs sont centrés et réduits; la constante n'est pas pénalisée. lambda=10 est choisi pour illustrer le calcul, sans prétendre être optimal. Une validation croisée est nécessaire pour le choisir selon la prévision.
\begin{Verbatim}
Z <- scale(as.matrix(donnees[variables]))
y <- donnees$mpg
lambda <- 10
beta_ridge <- solve(crossprod(Z) + lambda * diag(ncol(Z)),
                    crossprod(Z, y - mean(y)))
beta_ridge
pred_ridge <- drop(mean(y) + Z %*% beta_ridge)
\end{Verbatim}

\end{minipage}\par
\par\medskip\noindent\begin{minipage}{\linewidth}
\subsection*{Aide}
Dans R, utilisez \texttt{?lm} et l’aide des fonctions présentées. 
\end{minipage}\par
\end{document}
